As illustrated in Figure \ref{fig: embodied demo}, one application of \holodeck is synthesizing training environments to better match a novel testing distribution. To study this application, we consider ObjectNav~\cite{batra2020objectnav}, a common task in which a robot must navigate toward a specific object category. As existing benchmarks~\cite{deitke2020robothor, procthor,ramakrishnan2021hm3d} for ObjectNav consider only household environments and support a very limited collection of object types (16 object types in total combining the above benchmarks), we introduce \noveltythor, an artist-designed benchmark to evaluate embodied agents in diverse environments. Subsequently, we use the ObjectNav model pre-trained on \textsc{ProcTHOR-10K}~\cite{ai2thor} and finetune it on 100 scenes generated by \holodeck. These scenes are created by prompting \holodeck with the novel scene type as input. The model is then evaluated on \noveltythor.

\smallbreak
\noindent \textbf{\noveltythor.} We have two professional digital artists manually create 10 novel testing environments with two examples for each of the five categories: \textit{Office}, \textit{Daycare}, \textit{Music Room}, \textit{Gym}, and \textit{Arcade}. 
% We consider the task of ObjectNav (navigation towards a specific object category) in \noveltythor. 
Each scene contains novel object types not included in the existing ObjectNav tasks, e.g., ``piano'' in \textit{Music Room}, ``treadmill'' in \textit{Gym}, etc. Across \noveltythor, there are 92 unique object types.
\smallbreak

\noindent \textbf{Baselines.} For all methods except the one of random action, we use the same pre-trained ObjectNav model from \textsc{ProcTHOR-10K}~\cite{ai2thor}, which has been trained for ${\approx}$400M steps to navigate to 16 object categories. %As we are interested in understanding how \holodeck allows for adapting to new testing distributions, we finetune the above pre-trained model on 100 scenes generated by prompting \holodeck with the scene name as input.
To adapt the agent to novel scenes without human-construct training data, we consider two methods: (1) $+$\holodeck: we prompt\footnote{Here, we prompt with the scene name and its paraphrases to get more diverse outputs, e.g., we use ``game room'', ``amusement center'' for \textit{Arcade}.} \holodeck to generate 100 scenes for each scene type automatically; (2) $+$\textsc{Objaverse}: a strong baseline by enhancing \procthor with \holodeck's scene-type-specific object selection, specifically, those scenes are populated with similar Objaverse assets chosen by \holodeck.
% The main difference between $+$\textsc{Holodeck} and $+$\textsc{Objaverse} scenes is \holodeck's semantically meaningful layout design.

% Note that $+$\textsc{Objaverse} is an extremely strong baseline for ObjectNav as~\cite{maksymets2021treasurehunt} have found that navigating to \emph{arbirarily placed} objects already leads to strong ObjectNav models.


% use the same set of Objaverse assets by fine-tuning the same pre-trained model on scenes generated by \procthor enhanced by \holodeck's scene-type-specific object selection. That is, the generated scenes 

%\procthor to also select the Objaverse assets chosen by \holodeck for that scene type, i.e., we enhance \procthor with .

% More details about baselines and training can be found in the supplementary material.

\smallbreak
\noindent \textbf{Model.} Our ObjectNav models use the CLIP-based architectures of \cite{khandelwal2022:embodied-clip}, which contains a CNN visual encoder and a GRU to capture temporal information. We train each model with 100 scenes for 50M steps, which takes approximately one day on 8 Quadro RTX 8000 GPUs. We select the checkpoint of each model based on the best validation performance on its own validation scenes.
\smallbreak
\noindent \textbf{Results.} Table \ref{tab: embodied_results} shows zero-shot performance on \noveltythor. \holodeck achieves the best performance on average and surpasses baselines with considerable margins on \textit{Office}, \textit{Daycare}, and \textit{Music Room}. On \emph{Gym} and \emph{Arcade}, $+$\textsc{Holodeck} and $+$\textsc{Objaverse} perform similarly. Given that the main difference between $+$\textsc{Holodeck} and $+$\textsc{Objaverse} scenes is in the object placements, the observed difference suggests that \holodeck is more adept at creating layouts that resemble those designed by humans. For example, We can observe in Figure~\ref{fig: embodied demo} that the music room in \noveltythor contains a piano, violin cases, and cellos that are in close proximity to each other. The music room generated by \holodeck also shows a similar arrangement of these objects, highlighting the ``commonsense'' understanding of our method.
%crafted independently of \holodeck generated scenes by human experts, these novel environments This is also the case for the novel environments built by human experts. Note that these holodeck scenes are generated before those novel environemtns are built.
\procthor struggles in \noveltythor, often indistinguishably from random, because of poor object coverage during training.
%In summary, our experiments indicate that training with \holodeck can effectively aid agents in generalizing to new environments.

%\holodeck's meaningful design of layout.closer to 
%The is a strong baseline for ObjectNav as~\cite{maksymets2021treasurehunt} have found that navigating to \emph{arbirarily placed} objects already leads to strong ObjectNav models. 
%As noted before, our $+$\textsc{Objaverse} model is a strong baseline for the ObjectNav task, and we suspect \holodeck will have even greater value for other embodied tasks. 
% While it may seem surprising that the base \procthor model performs quite poorly across all environments, sometimes worse than Random, this is unavoidable as the pre-trained has never been trained on most object types in \noveltythor.


% It has not observed object targets during training and so can't navigate to them.
% \begin{figure}[!t]
% \centering
% \includegraphics[width=8.3cm]{images/object_type_performance.pdf}
%     \caption{}
%     \label{fig: single object performance}
%     \vspace{-0.5cm}
% \end{figure}
% Considering \procthor is already learned on 10K scenes, scaling cannot help the agent perform better on the novel scenes. 